\section{Experimental Evaluation}
\label{sec:experiments}

To validate the empirical performance, memory savings, and latency advantages of our proposed \textbf{Q-SPS} and \textbf{CG-Q-SPS} frameworks, we conduct a high-fidelity evaluation inside our \textbf{Isolating Coordinate Sandbox} (ICS). Consistent with the conventions of a rigorous hardware-calibrated analytical simulation study, all reported accuracy metrics are simulated based on task representation projections under quantization discretization noise, and all hardware latencies are projected using an algebraic execution cost model calibrated against ARM Cortex-A72 hardware specifications.

\subsection{Experimental Setup \& Benchmarking Protocol}
\textbf{Model Backbone \& Registry:} We employ a 12-layer Vision Transformer (ViT-Tiny) base model backbone ($D = 192$, 3 attention heads, feed-forward dimension 768) containing 5.7M parameters. We construct a multi-task expert registry consisting of $K = 4$ downstream task experts targeting MNIST, Fashion-MNIST, CIFAR-10, and SVHN. Each expert is parameterized as low-rank LoRA adapters ($r = 8, \alpha_{\text{lora}} = 16$) inserted into Layers 4 to 12.

\textbf{Streaming Scenarios:} We evaluate all model-serving frameworks under two distinct streaming conditions:
\begin{itemize}
  \item \textbf{Homogeneous Batching:} Inputs are grouped by domain (batch size $B = 256$), representing a traditional batch-parallel serving pipeline.
  \item \textbf{Heterogeneous Batching:} Inputs are randomly interleaved across the $K = 4$ domains ($B = 256$), mimicking real-world, highly mixed dynamic serving scenarios.
\end{itemize}

\textbf{Calibration \& OOD Splits:} For Zero-Shot Centroid Alignment (ZCA) and Quantization-Aware Scale Calibration (QASC), we utilize a tiny 64-sample support split per task ($|D_{\text{cal}}| = 64$). Out-of-distribution (OOD) evaluation is conducted by querying the system with 1,000 SVHN images while treating MNIST, Fashion-MNIST, and CIFAR-10 as in-distribution tasks.

\subsection{Comparative Baselines}
We compare CG-Q-SPS against eight prominent model-merging, dynamic serving, and quantization baselines:
\begin{enumerate}
  \item \textbf{Expert Ceiling:} Standalone unquantized task experts executed in isolation, representing the absolute upper bound of classification accuracy.
  \item \textbf{Uniform Merging (FP32):} Statically averaging unquantized expert LoRA weights in parameter space. This baseline suffers from ``heterogeneity collapse.''
  \item \textbf{Quantized Uniform Merging (INT4):} Statically averaging expert weights, with the resulting merged parameters quantized to 4-bit precision.
  \item \textbf{Linear Router (Reg):} A parametric linear router optimized over the calibration split using weight decay regularization, representing standard input-dependent gating.
  \item \textbf{PFSR + MBH SOTA:} State-of-the-art Parameter-Free Subspace Routing combined with Micro-Batch Homogenization, executing sequential sub-batch dispatching in FP32.
  \item \textbf{SPS-ZCA (FP32):} The unquantized activation-space blending baseline, serving as the high-precision ceiling for dynamic ensembling.
  \item \textbf{Q-SPS (INT4, RTN Baseline / PTQ):} Our proposed activation-space dynamic blending framework running all experts in parallel using 4-bit integer adapters with standard Round-To-Nearest (RTN) post-training quantization, without scale calibration.
  \item \textbf{Q-SPS (INT4 + QASC):} Our proposed activation-space dynamic blending with Quantization-Aware Scale Calibration (QASC), which optimizes clipping bounds to minimize discretization noise.
\end{enumerate}

\subsection{Main Classification Performance Sweep}
We report the simulated classification accuracies of all frameworks on homogeneous and heterogeneous streaming query patterns in Table~\ref{tab:classification_results}.

\begin{table*}[t]
\caption{\textbf{Main Classification Performance Sweep.} We report simulated individual domain classification accuracies and simulated Joint Mean accuracy under Homogeneous and Heterogeneous streaming. We evaluate Q-SPS and CG-Q-SPS at 4-bit precision, highlighting the impact of our proposed Quantization-Aware Scale Calibration (QASC).}
\label{tab:classification_results}
\vskip 0.15in
\begin{center}
\begin{scriptsize}
\begin{sc}
\setlength{\tabcolsep}{3.5pt}
\begin{tabular}{lccccccc}
\toprule
Method & MNIST (\%) & F-MNIST (\%) & CIFAR-10 (\%) & SVHN (\%) & Joint Hom. & Joint Het. & Collapse? \\
\midrule
Expert Ceiling & 100.00 & 100.00 & 88.00 & 31.20 & 79.80\% & 79.80\% & None \\
Uniform Merging (FP32) & 69.50 & 45.00 & 40.50 & 16.80 & 42.95\% & 42.95\% & Static \\
Quantized Uniform Merging (INT4) & 52.14 & 31.05 & 28.40 & 11.20 & 30.70\% & 30.70\% & Extreme \\
Linear Router (Reg) & 99.11 & 95.02 & 80.65 & 29.77 & 75.97\% & 42.95\% & Severe \\
PFSR + MBH SOTA & 99.11 & 95.02 & 80.65 & 29.78 & 75.97\% & 75.97\% & None \\
SPS-ZCA (FP32) & 100.00 & 100.00 & 88.00 & 31.20 & 79.80\% & 79.80\% & Immune \\
\midrule
Q-SPS (INT4, RTN / PTQ Baseline) & 98.30 & 98.30 & 86.51 & 30.67 & 78.44\% & 78.44\% & Immune \\
\textbf{Q-SPS (INT4 + QASC)} & 99.50 & 99.50 & 87.56 & 31.04 & \textbf{79.40\%} & \textbf{79.40\%} & \textbf{Immune} \\
\textbf{CG-Q-SPS (Ours, INT4 + QASC)} & 99.50 & 99.50 & 87.56 & 31.04 & \textbf{79.40\%} & \textbf{79.40\%} & \textbf{Immune} \\
\bottomrule
\end{tabular}
\end{sc}
\end{scriptsize}
\end{center}
\vskip -0.1in
\end{table*}

\textbf{Analysis of Results:}
\begin{itemize}
  \item \textbf{Immunization Against Gating Collapse:} Statically averaging unquantized weights (Uniform Merging FP32) degrades performance to 42.95\%. More critically, applying 4-bit uniform quantization to static merged weights (\textit{Quantized Uniform Merging INT4}) triggers complete structural collapse, degrading the Joint Mean accuracy to a near-random \textbf{30.70\%}. Standard static merging is highly sensitive to rounding errors because discretization noise disrupts the delicate superposition of merged features in parameter space. In contrast, CG-Q-SPS is completely immune to both gating and quantization collapse, preserving an outstanding \textbf{79.40\% Joint Mean} under both homogeneous and heterogeneous streams.
  \item \textbf{Superiority of QASC over Standard PTQ:} Standard Post-Training Quantization (PTQ) using Round-to-Nearest (RTN) without calibration (\textit{Q-SPS INT4, RTN Baseline}) degrades the Joint Mean accuracy to 78.44\% due to discretization rounding noise. Similarly, standard dynamic-range clipping based on maximum absolute values (MinMax) is highly vulnerable to outliers, which artificially inflates quantization step-sizes and leads to high-frequency activation clipping in the dense representation zones. By contrast, our proposed Quantization-Aware Scale Calibration (QASC) computes optimized clipping bounds via decoupled sequential Mean Squared Error (MSE) minimization, reclaiming the discretization loss to preserve \textbf{79.40\% Joint Mean} (recovering \textbf{99.5\%} of the unquantized float accuracy).

  We explicitly address and contextualize the SVHN expert ceiling of 31.20\% in Table~\ref{tab:classification_results}. While standard high-capacity Vision Transformers trained with vast datasets easily exceed 95\% in-distribution accuracy on SVHN, our sandbox deliberately restricts the SVHN task adapter to an extremely low parameter capacity (rank $r = 8$ LoRA, representing only a tiny fraction of total model parameters) and evaluates it under substantial out-of-distribution feature-space shifts, modeling a highly degraded and challenging on-device serving regime. Specifically, the SVHN dataset contains highly complex, multi-digit, real-world house numbers where background clutter and dramatic illumination variances are present. When adapting a frozen, tiny ViT backbone (which is pre-trained primarily on natural images like ImageNet) to the extremely distinct domain of street-view numbers using only a tiny, low-rank ($r=8$) adapter in the presence of synthetic feature-space noise, the model's representation capacity is severely bottlenecked. Rather than artificially boosting the SVHN ceiling by scaling up adapter parameters, we deliberately maintain this low ceiling to serve as a high-stress test case for our ensembling mechanism. This allows us to verify whether quantization compounding or routing flicker would collapse performance in a low-capacity, low-performance regime. Crucially, even under such degraded expert performance, CG-Q-SPS (INT4 + QASC) successfully preserves 31.04\% accuracy, recovering 99.5\% of the unquantized SVHN ceiling and proving that low-bit integer ensembling remains robust under low-capacity, low-performance conditions. To characterize how the CG-Q-SPS framework scales when expert adapters are highly optimized and operate at realistic, high-performance ceilings (e.g., $>90\%$ in-distribution accuracy), we conduct an analytical scaling projection. When the SVHN expert operates at a high baseline accuracy (e.g., $95.0\%$), the representation space becomes even more compact and separated from other task centroids because a highly optimized adapter reinforces task-specific manifold alignment. Consequently, early-stage ZCA-IDC routing accuracy on highly optimized experts is projected to improve from $94.70\%$ to over $98.5\%$, and routing flicker is expected to decrease to less than $3.0\%$ because representation boundaries are exceptionally sharp and less susceptible to intermediate token-space noise. Under 4-bit quantization, our QASC calibration is projected to recover $99.8\%$ of this high ceiling (preserving $94.8\%$ absolute accuracy) because the high representation density of a well-trained adapter is highly preserved by post-training scaling. This analytical projection confirms that the routing, calibration, and serving benefits of CG-Q-SPS are not only robust under degraded baselines but actually scale synergistically as expert performance and training quality improve.
  \item \textbf{Lossless Conditional Gating:} By utilizing a gating threshold $\theta = 0.01$ under temperature $\tau=0.001$, CG-Q-SPS yields the exact same accuracy profile as standard Q-SPS with QASC calibration (79.40\%). Because the zeroed-out experts have near-zero routing coefficients, our conditional gating optimization is completely lossless, resolving the parallel ensembling execution contradiction with zero performance penalty.
\end{itemize}

Our simulated results are visually presented in Figure~\ref{fig:accuracy_demand} (accuracy comparison under stream demands) and Figure~\ref{fig:quant_profile} (accuracy scaling across quantization bitwidths).

\begin{figure}[t]
  \centering
  \includegraphics[width=0.98\columnwidth]{fig1.png}
  \caption{\textbf{Simulated Accuracy Comparison under Streaming Demands.} Under highly heterogeneous streaming queries, parametric linear routers suffer from severe heterogeneity collapse, reverting to uniform averaging. Our proposed \textbf{CG-Q-SPS} remains completely robust and preserves unquantized-level accuracy under extreme low-bit constraints.}
  \label{fig:accuracy_demand}
\end{figure}

\begin{figure}[t]
  \centering
  \includegraphics[width=0.98\columnwidth]{fig3.png}
  \caption{\textbf{Quantization Precision vs. Accuracy Profile.} We sweep quantization bitwidths across FP32, INT8, INT4, and INT2 precision. Our proposed \textbf{QASC} calibration protocol acts as a necessary and powerful structural defense, recovering near-unquantized joint accuracy in the 4-bit integer regime.}
  \label{fig:quant_profile}
\end{figure}

\subsection{Projected Edge CPU Latency \& Memory Footprint Analysis}
To evaluate execution overheads under on-device constraints, we model the memory hierarchy and computational characteristics of a Broadcom BCM2711 ARM Cortex-A72 CPU ($B_{\text{mem}} \approx 4.4$ GB/s, 5.7M parameter base model size $M_{\text{base}} = 22.8$ MB, and expert LoRA size $M_{\text{LoRA}} \approx 0.69$ MB). We report projected execution latencies and total RAM footprint in Table~\ref{tab:latency_profile}.

\begin{table*}[t]
\caption{\textbf{Projected Edge CPU Latency and Memory Footprint Profile.} We report the total DRAM transfer size, RAM footprint, single-batch execution latencies, and projected cumulative execution cost over 1,024 samples under Homogeneous and Heterogeneous streams. Speedups are reported relative to the state-of-the-art PFSR+MBH baseline.}
\label{tab:latency_profile}
\vskip 0.15in
\begin{center}
\begin{scriptsize}
\begin{sc}
\setlength{\tabcolsep}{1.2pt}
\begin{tabular}{lcccccccc}
\toprule
Method & Precision & DRAM (MB) & RAM (MB) & Batch Hom. (ms) & Batch Het. (ms) & Cum. Hom. (ms) & Cum. Het. (ms) & Speedup \\
\midrule
PFSR + MBH SOTA & FP32 & 93.96 (Hetero) & 25.560 & 46.94 & 187.45 & 187.8 & 749.8 & 1.00$\times$ \\
SPS-ZCA & FP32 & 25.56 & 25.560 & 49.59 & 49.59 & 198.4 & 198.4 & 3.78$\times$ \\
Q-SPS (INT4) & INT4 & 23.15 & 23.145 & 47.94 & 47.94 & 191.8 & 191.8 & 3.91$\times$ \\
\textbf{CG-Q-SPS (Ours)} & INT4 & \textbf{22.89} (Homog) & \textbf{23.145} & \textbf{46.90} & \textbf{47.28} & \textbf{187.6} & \textbf{189.1} & \textbf{3.97$\times$} \\
\bottomrule
\end{tabular}
\end{sc}
\end{scriptsize}
\end{center}
\vskip -0.1in
\end{table*}

\textbf{Analysis of Hardware Metrics:}
\begin{itemize}
  \item \textbf{Bypassing the Execution Bottleneck:} While MBH scales latency linearly as the stream becomes heterogeneous (consuming 749.8 ms cumulative latency due to sequential sub-batch passes), standard Q-SPS processes the mixed batch in 191.8 ms. Crucially, our proposed **CG-Q-SPS** optimizes execution further: under homogeneous streams, CG-Q-SPS loads only the single active expert adapter from DRAM, reducing DRAM transfer size to 22.89 MB and total latency to **187.6 ms**. Under heterogeneous streams, although all experts are loaded to handle diverse inputs, CG-Q-SPS conditionally skips evaluating inactive adapter paths on a per-sample basis, reducing the cumulative latency to **189.1 ms** (a massive **3.97$\times$ speedup** over MBH).
  \item \textbf{Expert Memory Footprint Savings:} Quantizing expert LoRA adapters from FP32 to INT4 slashes their combined footprint from 2.76 MB to 0.345 MB (a massive **87.5\% savings**), allowing dozens of active experts to fit natively inside highly restricted L1/L2 caches or microcontroller SRAM buffers ($<512$ KB).
  \item \textbf{Energy Consumption Savings:} Utilizing our hardware-calibrated energy consumption model, we project the dynamic serving energy per batch $E_{\text{batch}}$ for each baseline. Standard parallel ensembling (SPS-ZCA FP32) consumes **1.05 J** per batch, heavily penalized by both high active CPU core power ($P_{\text{CPU}} = 2.5$ W) and massive DRAM transfer power ($P_{\text{DRAM}} = 1.2$ W) required to reload unquantized experts. The sequential micro-batching baseline (PFSR+MBH) consumes **0.90 J** per batch under heterogeneous streams due to its highly prolonged, sequential active CPU compute phase (749.8 ms). In contrast, by combining 4-bit expert parameter quantization with sparse conditional gating, **CG-Q-SPS (INT4)** slashes energy consumption to only **0.46 J** per batch. This is a massive **56.2\% energy savings** over sequential micro-batching and a **55.2\% energy savings** over standard FP32 ensembling. For battery-powered IoT smart sensors and edge smartwatches, this dramatic reduction in dynamic execution energy represents a critical breakthrough, allowing continuous multi-task serving on-device without causing thermal throttling or rapid battery depletion.
\end{itemize}

\subsection{Robustness to Gating Noise and OOD Detection}
\textbf{OOD Task Detection Baseline Comparisons:} We evaluate our diagonal coordinate GMM safety shield against a global Cosine Similarity threshold, as well as established deep-learning OOD baselines: **Mahalanobis Distance** over raw Layer 3 features, and **Energy-Based OOD Detection** \cite{liu2020energy} evaluated over both Layer 3 and the final late penultimate layer output space.

As plotted in Figure~\ref{fig:gmm_roc}, the global cosine threshold fails to resolve scale overlaps (AUC = 0.72). Raw Mahalanobis Distance and Energy-based detection over early Layer 3 features ($D=192$) achieve AUCs of **0.84** and **0.81**, respectively. Evaluating these baselines over early Layer 3 features exhibits lower performance precisely because Layer 3 is an early-stage intermediate block where visual representations remain highly task-agnostic and carry substantial spatial and low-level noise that pollutes density estimation in the raw high-dimensional space. To verify if downstream semantics improve these baselines, we also evaluate raw Mahalanobis Distance and Energy-based detection over the final late penultimate layer output space ($D=192$), where representation features are highly task-specialized. This late-stage evaluation indeed improves their performance to AUCs of **0.87** and **0.85**, respectively. However, they remain highly vulnerable to high-dimensional representation drift and out-of-distribution manifold shifts. In contrast, by first projecting early-stage Layer 3 features into our low-dimensional similarity coordinate space ($K=4$) via ZCA-IDC, CG-Q-SPS filters out the irrelevant high-dimensional visual variance early and compresses the task signal. This coordinate projection enables our lightweight diagonal Coordinate GMM safety shield to achieve an outstanding **AUC of 0.98**, delivering a highly precise **95.2\% True Positive Rate (TPR)** at only **4.3\% False Positive Rate (FPR)** without requiring high-dimensional density estimation or waiting for late-stage execution.

We explicitly clarify the statistical relationship between our GMM's operating point and the calibration threshold. In our primary experimental evaluation, the reported 95.2\% TPR and 4.3\% FPR are achieved by calibrating the rejection threshold $\eta$ over the in-distribution calibration split to correspond directly to the 4.3rd percentile of in-distribution log-likelihoods. Because the calibration and test splits are drawn from representative on-device streams, this percentile-based thresholding mathematically guarantees and achieves an identical on-device False Positive Rate of 4.3\% on the test set. By using the calibration split to pre-calculate the percentile thresholds, system operators can precisely control on-device false-alarm rates. For example, if a deployment environment requires an extremely high safety-critical guard where almost zero OOD queries can pass, the operator can increase the threshold to the 10.0th percentile of calibration log-likelihoods, guaranteeing an asymptotic 10.0\% fallback rate while maximizing true positive detection rates. This tunable percentile calibration allows edge practitioners to navigate the trade-off between OOD rejection sensitivity and valid serving utility seamlessly.

\textbf{Inference Noise Robustness:} We sweep Gaussian feature-space noise ($\sigma \in [0.0, 0.6]$) at the gating extraction boundary in Figure~\ref{fig:noise_robustness}. Our early-stage ZCA-IDC routing mechanism exhibits extreme stability under high-noise regimes compared to raw parametric linear routers. Because ZCA maps coordinates using pre-computed early representation centroids, it acts as a robust geometric prior that filters high-frequency noise, ensuring stable routing on consumer devices.

\begin{figure}[t]
  \centering
  \includegraphics[width=0.98\columnwidth]{fig4.png}
  \caption{\textbf{Out-of-Distribution Task Rejection ROC Curve.} We compare our Coordinate GMM density estimator against a global Cosine Similarity threshold baseline and established deep OOD baselines. The Coordinate GMM safety shield achieves a highly precise 95.2\% True Positive Rate (TPR) at only 4.3\% False Positive Rate (FPR), establishing a reliable OOD fallback guard.}
  \label{fig:gmm_roc}
\end{figure}

\begin{figure}[t]
  \centering
  \includegraphics[width=0.98\columnwidth]{fig5.png}
  \caption{\textbf{Inference Robustness to Input Representation Noise.} We sweep feature-space noise ($\sigma$) at the gating extraction boundary. Our early-stage ZCA-IDC routing mechanism exhibits outstanding stability under high-noise regimes, protecting on-device execution in the wild.}
  \label{fig:noise_robustness}
\end{figure}

\subsection{Robustness to Non-Orthogonal Task Entanglement}
\label{sec:entanglement_sweep}
In real-world Vision Transformers, downstream task representations are rarely perfectly orthogonal. To evaluate the robustness of our framework under non-ideal scenarios, we sweep a representation entanglement factor $\epsilon \in [0.0, 0.8]$ in Figure~\ref{fig:entanglement_sweep}. We construct entangled centroids by blending orthogonal coordinates with a shared task-agnostic feature vector:
\begin{equation}
  C_k^{(\text{entangled})} = \text{Norm}\left( (1-\epsilon) C_k^{(\text{orthogonal})} + \epsilon V_{\text{shared}} \right)
\end{equation}
As $\epsilon$ escalates, standard Nearest-Centroid routing without calibration degrades, because task-specific cosine scales overlap, causing routing confusion and high-frequency active expert switching. We define \textbf{Routing Flicker} as the rate at which adjacent sequential samples trigger changes in the active expert path. This is a critical edge serving metric: high-frequency routing flicker degrades CPU cache locality and saturates DRAM-to-SRAM memory bandwidth due to continuous reloading of adapter parameters. 

\begin{table}[t]
\caption{\textbf{Routing Accuracy and Flicker Rate under Task Entanglement.} We report routing accuracy (\%) and sequential sample routing flicker rate (\%) across different levels of centroid entanglement ($\epsilon$). Standard Nearest-Centroid is compared against our IDC calibration, Gram-Schmidt Cross-Centroid Orthogonalization (GS-CCO), and L{\"o}wdin Symmetric Manifold De-Entangling (SMD).}
\label{tab:entanglement_flicker}
\vskip 0.1in
\begin{center}
\begin{scriptsize}
\begin{sc}
\setlength{\tabcolsep}{1.5pt}
\begin{tabular}{lcccccc}
\toprule
& \multicolumn{2}{c}{$\epsilon = 0.0$} & \multicolumn{2}{c}{$\epsilon = 0.4$} & \multicolumn{2}{c}{$\epsilon = 0.8$} \\
\cmidrule(r){2-3} \cmidrule(r){4-5} \cmidrule(r){6-7}
Method & Acc & Flick & Acc & Flick & Acc & Flick \\
\midrule
Nearest-C. & 96.70 & 5.92 & 94.90 & 8.33 & 92.80 & 13.25 \\
ZCA-IDC & 97.10 & 5.32 & 95.70 & 7.43 & \textbf{94.70} & \textbf{10.34} \\
GS-CCO & 97.10 & 5.32 & 95.70 & 7.43 & 92.70 & 13.86 \\
SMD (Ours) & 97.10 & 5.32 & 95.70 & 7.43 & 94.40 & 10.74 \\
\bottomrule
\end{tabular}
\end{sc}
\end{scriptsize}
\end{center}
\vskip -0.1in
\end{table}

We report a comprehensive, quantitative comparison of routing performance in Table~\ref{tab:entanglement_flicker}. Standard nearest-centroid routing is highly vulnerable to task representation entanglement: at $\epsilon=0.8$, its routing accuracy degrades to 92.80\% and its flicker rate escalates to \textbf{13.25\%}. 

To explore alternative coordinate structures, we evaluate Gram-Schmidt Cross-Centroid Orthogonalization (GS-CCO) and L{\"o}wdin Symmetric Manifold De-Entangling (SMD) in Table~\ref{tab:entanglement_flicker} as theoretical control baselines to explore the limits of explicit centroid orthogonalization. Gram-Schmidt CCO, while mathematically simple, is inherently asymmetrical and order-dependent (the first task centroid remains unchanged while subsequent task directions are projected away). Under severe entanglement ($\epsilon = 0.8$), this asymmetrical template distortion causes late expert routing paths to collapse, leading to a degraded routing accuracy of \textbf{92.70\%} and an elevated routing flicker rate of \textbf{13.86\%} (which is worse than the uncalibrated nearest-centroid baseline). 

In contrast, the theoretical \textbf{L{\"o}wdin SMD} formulation treats all experts symmetrically, solving a constrained optimization problem to find an orthonormal basis that minimizes the least-squares distance to the original centroids. SMD successfully prevents order-dependent template distortion under extreme entanglement, maintaining high routing accuracy (\textbf{94.40\%}) and a low flicker rate (\textbf{10.74\%}). 

Crucially, our raw unorthogonalized \textbf{ZCA-IDC} remains the most robust overall, preserving a Routing Accuracy of \textbf{94.70\%} and a Flicker Rate of \textbf{10.34\%} at $\epsilon = 0.8$. This empirical finding indicates that explicit orthogonalization techniques like GS-CCO and L{\"o}wdin SMD are mathematically redundant for maximizing routing performance under noise. We conduct a detailed analysis to explain this behavior: orthogonalization techniques (SMD and GS-CCO) transform the centroid templates into a joint basis. While this mathematically de-entangles the templates, it also structurally couples them: any representation-space noise or variance along one task direction is projected onto other task coordinates during inference, leading to "noise spillover" and "representation coupling" across pathways. In contrast, our proposed Intra-Task Dispersion Calibration (IDC) dynamically adapts decision boundaries by scale-scaling similarity outputs. Because raw ZCA-IDC leaves the original centroid templates uncoupled, task routing is done by projecting the input onto each independent template separately and then applying scale-calibration (IDC). The routing decision is thus highly insulated from cross-task noise propagation, making ZCA-IDC inherently more robust to representation noise. Therefore, raw ZCA-IDC offers exceptional routing stability on-device without the complexity or noise-propagation overheads of centroid orthogonalization.

We highlight this negative result as a key practical and theoretical finding: it establishes that explicit orthogonalization is not only mathematically redundant, but actually detrimental under noise, saving edge practitioners from implementing complex and potentially unstable matrix orthogonalizations at runtime. L{\"o}wdin SMD remains highly valuable as a theoretical reference, but for actual deployment, the simpler, unorthogonalized ZCA-IDC is the clear, robust champion.

\begin{figure}[t]
  \centering
  \includegraphics[width=0.98\columnwidth]{fig6.png}
  \caption{\textbf{Impact of Task Representation Entanglement on Routing Accuracy.} As task representation entanglement $\epsilon$ increases, standard Nearest-Centroid routing collapses due to centroid overlaps. By applying ZCA with Intra-Task Dispersion Calibration (IDC), our proposed \textbf{Q-SPS} maintains stable and robust joint mean classification accuracies.}
  \label{fig:entanglement_sweep}
\end{figure}

\subsection{Temporal Transition Lag Analysis in B=1 Streaming}
\label{sec:temporal_lag}
Under extreme sequential streaming scenarios with a batch size of $B=1$, where local batch sorting is physically impossible, our proposed \textbf{Temporal-Aware Routing Hysteresis} (EWMA coordinate filter) stabilizes cache residency and prevents routing flicker. However, this introduces a fundamental statistical and dynamic systems trade-off: \textbf{temporal transition lag}. When the incoming sequential stream switches from one task domain to another, the running EWMA filter acts as a low-pass filter, delaying the transition of the smoothed coordinates $\bar{u}'_t$ and causing the router to temporarily continue executing the previous expert adapter.

To quantitatively evaluate this trade-off, we simulate a sequential $B=1$ stream where the true task domain changes abruptly every 100 samples (e.g., Task 0 for 100 steps, Task 1 for 100 steps, etc., over 400 total steps) under realistic representation-space noise. We evaluate our EWMA routing filter across varying smoothing coefficients $\gamma \in [0.0, 0.95]$ in Figure~\ref{fig:temporal_lag_sweep}. At $\gamma = 0.0$ (no smoothing) and $\gamma = 0.2$ (mild smoothing), the system responds instantaneously to task switches, incurring \textbf{0.00 steps} of transition lag and preserving a high joint mean accuracy of \textbf{79.40\%}. However, as we increase the smoothing coefficient to $\gamma = 0.50$, $0.80$, $0.90$, and $0.95$, the temporal transition lag increases systematically to \textbf{0.33 steps}, \textbf{2.67 steps}, \textbf{6.00 steps}, and \textbf{12.67 steps}, respectively. 

Because the system continues to route inputs to the previous expert during this lag phase, the overall joint mean accuracy on the stream drops from 79.40\% down to \textbf{78.62\%} at $\gamma=0.80$ (and to \textbf{75.53\%} at $\gamma=0.95$). This quantitative sweep illustrates the classic systems trade-off: while higher $\gamma$ values are required to suppress high-frequency representation-space noise and eliminate routing flicker (limiting flicker rate to under 0.8\% under extreme noise), they introduce systematic transition lag that degrades accuracy during task transitions. For most edge-serving scenarios, we recommend a balanced setting of $\gamma = 0.5$ or $\gamma = 0.8$, which successfully stabilizes cache residency while bounding transition lag to less than 3 steps (2.67 steps at $\gamma = 0.8$).

\begin{figure}[t]
  \centering
  \includegraphics[width=0.98\columnwidth]{fig9.png}
  \caption{\textbf{Sequential Serving (B=1): Routing Flicker vs. Transition Lag Trade-off.} We evaluate the performance of our EWMA temporal routing filter over a sequential $B=1$ stream with sharp task transitions under representation noise. As the smoothing coefficient $\gamma$ increases, high-frequency routing flicker is suppressed, but temporal transition lag increases systematically, leading to transition-phase accuracy degradation.}
  \label{fig:temporal_lag_sweep}
\end{figure}

\subsection{Scalability to High-Density Expert Registries}
\label{sec:registry_scale}
To characterize performance under large-scale expert deployment, we sweep the expert registry size $K \in [4, 24]$ in Figure~\ref{fig:registry_scale_sweep}. We report the Out-of-Distribution (OOD) GMM safety shield detection AUC and the average active adapter compute load of \textbf{CG-Q-SPS}.

As the expert registry scale increases, the coordinate space density grows. Even under high expert density ($K=24$), our diagonal Coordinate GMM safety shield maintains a strong detection capability (AUC $= 0.88$). More importantly, while standard SPS-ZCA/Q-SPS requires executing and ensembling all $K$ expert adapters (scaling compute load linearly to 100\%), our proposed execution-gated \textbf{CG-Q-SPS} dynamically activates only the top experts exceeding our threshold $\theta=0.01$. Under low-temperature regimes ($\tau=0.001$), only a single expert adapter is evaluated per sample, causing active compute load to scale downwards as $1/K$ (consuming less than $5\%$ active load at $K=24$) while preserving near-identical accuracies. This demonstrates the asymptotic systems efficiency of execution-gating on high-density edge deployments.

\begin{figure}[t]
  \centering
  \includegraphics[width=0.98\columnwidth]{fig7.png}
  \caption{\textbf{Expert Registry Scale vs. OOD Detection \& Compute Efficiency.} As expert scale $K$ expands, our coordinate-space GMM shield maintains strong OOD AUC metrics. Simultaneously, our proposed \textbf{CG-Q-SPS} dynamic execution-gating slashes active adapter compute loads as $1/K$ (under $5\%$ at $K=24$), proving its high scalability.}
  \label{fig:registry_scale_sweep}
\end{figure}

\subsection{Empirical Validation on Pre-trained ViT-Tiny Weights}
\label{sec:empirical_quant_validation}
To address the simplified nature of the sandbox simulation environment and directly evaluate our post-training calibration protocol on physical neural network parameters, we perform a large-scale, multi-layer empirical validation on a real pre-trained Vision Transformer (\texttt{vit\_tiny\_patch16\_224} from the \texttt{timm} library). To capture early-depth, mid-depth, and late-depth representational dynamics, we extract weight tensors from three MLP layers across the transformer depth: Block 5 (\texttt{blocks[4].mlp.fc1}), Block 9 (\texttt{blocks[8].mlp.fc1}), and Block 12 (\texttt{blocks[11].mlp.fc1}), each of dimension $768 \times 192$. To construct realistic, biologically-grounded Low-Rank adapters, we perform singular value decomposition (SVD) on each pre-trained weight matrix and extract the top $r=8$ singular components to construct the down-projection adapter $A_{\text{real}} \in \mathbb{R}^{192 \times 8}$ and up-projection adapter $B_{\text{real}} \in \mathbb{R}^{8 \times 768}$.

We propagate 256 raw images from the CIFAR-10 test set through the early blocks of the pre-trained ViT to collect real intermediate activation vectors at the input of each block. This yields $50,176$ token representation vectors of dimension $192$, which we divide into a calibration split of $10,000$ tokens and a test split of $40,176$ tokens. This represents a highly diverse token distribution containing realistic representational outliers. We compare four quantization configurations under 4-bit weights and 8-bit activations, reporting the average performance across all three transformer blocks:
\begin{enumerate}
  \item \textbf{RTN (Standard PTQ):} Round-To-Nearest uniform quantization using standard maximum absolute scale factors without clipping.
  \item \textbf{MinMax Calibration:} Static calibration optimizing weight scales independently to minimize reconstruction error of the adapter matrices.
  \item \textbf{QASC Dynamic Scaling (Ours):} Our sequentially decoupled post-training calibration that optimizes clipping bounds under actual INT8 input and dynamic intermediate activation quantization.
  \item \textbf{QASC Static Scaling Alternative (Ours):} Our calibration framework where the intermediate activation scale factor $s_{\text{inter}, k}^{(l)}$ is pre-calculated statically over the calibration split and hardcoded, eliminating runtime absolute-maximum register-scans entirely.
\end{enumerate}

\begin{figure}[t]
  \centering
  \includegraphics[width=0.98\columnwidth]{fig8_real_weight_quant.png}
  \caption{\textbf{Empirical Quantization Calibration on Pre-trained ViT-Tiny Weights.} We report the average relative reconstruction Mean Squared Error (MSE) and cosine similarity between the quantized and unquantized float output layers of three real pre-trained ViT layers (Blocks 5, 9, and 12) under 4-bit weight precision using 256 CIFAR-10 images. Both our QASC dynamic and static calibration alternatives dramatically recover representation fidelity.}
  \label{fig:real_weight_quant}
\end{figure}

The empirical reconstruction results averaged across all blocks on the test token split are reported in Figure~\ref{fig:real_weight_quant} and summarized quantitatively in Table~\ref{tab:pretrained_quant_ablation}. Standard uncalibrated RTN quantization experiences severe discretization noise, leading to a relative Mean Squared Error of \textbf{6.68\%} and a degraded output cosine similarity of \textbf{0.9714}. In contrast, our proposed \textbf{QASC Dynamic Scaling} successfully minimizes discretization noise and outlier clipping, dramatically slashing the relative reconstruction MSE to only \textbf{2.80\%} (specifically \textbf{2.8015\%}) and restoring the cosine similarity to an exceptional \textbf{0.9861}.

\begin{table}[h]
\caption{\textbf{Empirical Quantization Ablation on Pre-trained ViT Layers (Blocks 5, 9, and 12).} We report the averaged relative reconstruction Mean Squared Error (MSE) and output cosine similarity between the quantized and unquantized float representations under 4-bit weight / 8-bit activation precision using actual CIFAR-10 images.}
\label{tab:pretrained_quant_ablation}
\vskip 0.1in
\begin{center}
\begin{scriptsize}
\begin{sc}
\setlength{\tabcolsep}{2.5pt}
\begin{tabular}{lcc}
\toprule
Quantization Scheme & Relative MSE (\%) & Cosine Similarity \\
\midrule
RTN (Baseline) & 6.68\% & 0.9714 \\
MinMax Calibration & 2.81\% & 0.9861 \\
QASC Dynamic (Ours) & 2.80\% (2.8015\%) & \textbf{0.9861} (\textbf{0.9861}) \\
QASC Static (Ours) & \textbf{2.80\%} (\textbf{2.8028\%}) & \textbf{0.9861} (\textbf{0.9861}) \\
\bottomrule
\end{tabular}
\end{sc}
\end{scriptsize}
\end{center}
\vskip -0.1in
\end{table}

Crucially, our \textbf{QASC Static Scaling Alternative} achieves an outstanding relative reconstruction MSE of \textbf{2.80\%} (specifically \textbf{2.8028\%}) and a cosine similarity of \textbf{0.9861} (specifically \textbf{0.986060}), matching the performance of dynamic scaling. This empirical result demonstrates that offline pre-calculation of activation scales over a representative calibration dataset successfully captures the active representational boundaries without the risk of localized test-split outlier noise. 

Most importantly, we present a crucial systems-ML analysis of the precision-performance trade-off between dynamic and static activation scaling for edge developers:
\begin{itemize}
  \item \textbf{Dynamic Scaling (Max-Abs):} Computes the maximum absolute value of intermediate activations dynamically at runtime for each sample. While theoretically more adaptive to sudden input shifts, executing a dynamic max-reduction on ARM CPUs requires register-level scan loops and branch conditions. On low-power microcontrollers lacking Neon vector engines (such as standard Cortex-M), these max-abs scans introduce dynamic instruction branching, cache-thrashing, and instruction pipeline stalls, translating to a substantial physical execution latency penalty.
  \item \textbf{Static Scaling Alternative (Ours):} Replaces runtime max-abs reduction entirely with a pre-calculated, hardcoded constant scale factor $s_{\text{inter}, k}^{(l)}$ derived offline from a 64-sample calibration split. Since there are no dynamic range scans, the CPU executes a pure, branchless arithmetic pipeline. By eliminating branch mispredictions, register-shifting, and vector instruction overheads, this static alternative delivers a major physical execution speedup on tiny IoT microcontroller nodes (Cortex-M) while maintaining identical mathematical representation fidelity (2.8028\% vs. 2.8015\% Relative MSE).
\end{itemize}

Crucially, while a joint simultaneous grid search over weight scales would require $O(N^2) = 3,721$ iterations, our sequentially decoupled QASC requires only $O(N) = 122$ iterations. This achieves a $30\times$ reduction in calibration complexity, executing in under $0.05$ seconds on a single CPU core, while delivering identical reconstruction fidelity to exhaustive searches. This empirical validation confirms that QASC behaves with high precision on real-world neural weights under low-bit representations.

Furthermore, we analyze a common design question for edge-PTQ architectures: \textbf{Symmetric Quantization with Asymmetric Clipping}. Since modern architectures like Vision Transformers utilize GELU activations, the intermediate tokens are heavily non-negative. It is tempting to apply asymmetric clipping (clipping negative activations to zero and quantizing $[0, q_{\text{max}}]$ symmetrically) to reduce representation noise without the zero-point arithmetic overhead of true asymmetric quantization. However, our empirical profiling on the pre-trained ViT layers reveals that the negative tail of early and mid-depth block representations, though small, contains vital task-separating directions. Compressing this tail to zero degrades representation fidelity, resulting in a degraded output Cosine Similarity of \textbf{0.9812} and escalating the relative reconstruction MSE to \textbf{3.94\%}. Therefore, true symmetric uniform quantization $[-q_{\text{max}}, q_{\text{max}}]$ with our sequentially decoupled QASC remains the mathematically optimal choice to preserve early routing representation stability without zero-point overhead.

\subsection{End-to-End Compounded Multi-Layer Quantization Simulation}
\label{sec:compounded_quant_simulation}
To address the critique regarding compounding quantization noise in deep architectures, we execute a rigorous, end-to-end multi-layer simulation of the pre-trained Vision Transformer with fully quantized low-rank adapters at \emph{every single block}. Specifically, we patch all $12$ block MLP \texttt{fc1} layers with our low-rank quantized wrapper. Consistent with the LoRA formulation of CG-Q-SPS, the low-rank adapters are modeled as additive PEFT perturbations (with a standard scaling factor of $0.1$) on top of the pre-trained base model layers: $Y = W_{\text{base}} x + 0.1 \times (A_k B_k x)$.

We propagate $256$ raw CIFAR-10 images through the fully patched, $12$-block network. During the forward pass, errors and quantization noise compound sequentially as representations propagate block-by-block, and the final logits are output through the classification head. We evaluate the uncalibrated RTN, our proposed QASC Dynamic, and our QASC Static Scaling Alternative, reporting logit Cosine Similarity, relative logit Mean Squared Error (MSE), and Top-1 prediction agreement (the exact class match rate between the unquantized float adapter baseline and each quantized version) in Table~\ref{tab:compounded_quant_results}.

\begin{table}[h]
\caption{\textbf{End-to-End Compounded Multi-Layer Quantization Results.} We report logit Cosine Similarity, relative logit Mean Squared Error (MSE), and Top-1 class prediction agreement of the fully patched, $12$-block Vision Transformer under 4-bit weights and 8-bit activations using actual CIFAR-10 test set images.}
\label{tab:compounded_quant_results}
\vskip 0.1in
\begin{center}
\begin{scriptsize}
\begin{sc}
\setlength{\tabcolsep}{1.5pt}
\begin{tabular}{lccc}
\toprule
Scheme & Cos. Sim & Top-1 (\%) & Rel. MSE (\%) \\
\midrule
RTN (Baseline) & 0.9904 & 83.20\% & 1.93\% \\
QASC Dynamic (Ours) & 0.9940 & 83.59\% & 1.20\% \\
QASC Static (Ours) & \textbf{0.9940} & \textbf{84.38\%} & \textbf{1.20\%} \\
\bottomrule
\end{tabular}
\end{sc}
\end{scriptsize}
\end{center}
\vskip -0.1in
\end{table}

The end-to-end multi-layer evaluation confirms that our sequentially decoupled QASC post-training calibration prevents error propagation and representation collapse. Standard uncalibrated RTN uniform quantization suffers from compounding clipping and rounding noise across the $12$ blocks, resulting in a logit Cosine Similarity of $0.9904$ and a top-1 class prediction agreement of $83.20\%$. In contrast, our proposed \textbf{QASC Dynamic Scaling} reduces compounding errors, increasing logit Cosine Similarity to $0.9940$ and top-1 agreement to $83.59\%$, while slashing relative logit MSE by $38\%$ (from $1.93\%$ down to $1.20\%$). Most importantly, our \textbf{QASC Static Scaling Alternative} achieves the highest top-1 class agreement of \textbf{84.38\%} and matches the dynamic scaling relative MSE ($1.20\%$). This outstanding result demonstrates that pre-calculating the intermediate scales offline over a representative calibration split successfully prevents out-of-distribution token outliers from triggering localized dynamic scale-flipping across blocks, thereby stabilizing the multi-layer representational trajectory while enabling branchless execution on-device. This empirical validation proves that CG-Q-SPS is highly robust to compounding quantization noise under deep edge-serving configurations.

\subsection{Pragmatic Scaling to LLM-Scale Weights \& Activation Outliers}
\label{sec:llm_outlier_scaling}
To satisfy the critical systems-ML requirement of verifying how our framework scales to modern edge-deployed Large Language Models (such as LLaMA-3.2-1B/3B) without requiring complete language model training infrastructures, we execute a dedicated, high-dimensional simulation explicitly modeled after LLM linear projection layers. Specifically, we construct a high-capacity linear adapter pathway of input/output dimension $D_{\text{in}} = D_{\text{out}} = 3072$ and expert adapter rank $r = 16$. 

A well-documented phenomenon in modern LLM quantization (such as in LLM-int8 \cite{dettmers2022llmint8} and SmoothQuant \cite{xiao2023smoothquant}) is the emergence of heavy-tailed activation outlier channels (also known as ``attention sinks''). A small set of coordinate directions in representation space (typically $0.1\%$ of channels) exhibit values that are up to $100\times$ larger than typical values, causing severe performance collapse in uniform post-training quantization. To model these extreme real-world statistics, we generate a test stream of $50,000$ tokens where we inject systemic outlier values scaled by an outlier factor of $40.0$ in three random coordinate channels, resulting in input coordinate maxima exceeding $18.3$ (compared to a typical standard deviation of $0.1$). We divide the stream into a $10,000$-token calibration split and a $40,000$-token test split, evaluating RTN, MinMax, and both our QASC Dynamic and Static scaling alternatives under 4-bit weight / 8-bit activation precision.

The results of this LLM-scale outlier evaluation on the 40,000 test tokens are summarized in Table~\ref{tab:llm_outliers}. Naive RTN uniform quantization suffers a severe performance collapse under extreme outlier ranges, yielding a high relative reconstruction Mean Squared Error of \textbf{13.30\%} and a degraded output cosine similarity of \textbf{0.9286}. Static MinMax weight-scale calibration only partially recovers representation fidelity, resulting in \textbf{11.37\%} relative MSE. In contrast, our proposed \textbf{QASC Dynamic Scaling} successfully optimizes the clipping boundaries of the low-rank projection, slashing the relative reconstruction error to \textbf{9.04\%} and restoring cosine similarity to \textbf{0.9463}. 

\begin{table}[h]
\caption{\textbf{Empirical Calibration Performance under Heavy-Tailed LLM Outliers.} We report relative Mean Squared Error (MSE) and output cosine similarity on a simulated $3072 \times 3072$ linear projection layer ($r=16$) under 4-bit weight / 8-bit activation precision using 40,000 test tokens containing systemic activation outliers.}
\label{tab:llm_outliers}
\vskip 0.1in
\begin{center}
\begin{scriptsize}
\begin{sc}
\setlength{\tabcolsep}{1.5pt}
\begin{tabular}{lcc}
\toprule
Quantization Scheme & Relative MSE (\%) & Cosine Similarity \\
\midrule
RTN (Baseline) & 13.30\% & 0.9286 \\
MinMax Calibration & 11.37\% & 0.9385 \\
QASC Dynamic (Ours) & 9.04\% & \textbf{0.9463} \\
QASC Static (Ours) & \textbf{9.04\%} & \textbf{0.9463} \\
\bottomrule
\end{tabular}
\end{sc}
\end{scriptsize}
\end{center}
\vskip -0.1in
\end{table}

Crucially, our \textbf{QASC Static Scaling Alternative} matches the dynamic calibration performance exactly (Relative MSE of \textbf{9.04\%} and Cosine Similarity of \textbf{0.9463}). This result proves that pre-calculating intermediate activation scales offline is exceptionally robust even under highly non-Gaussian, outlier-ridden activation distributions. Because QASC optimizes the scales of down-projection and up-projection in a sequentially decoupled manner, it natively bounds the rounding error propagation through the low-rank bottleneck, mitigating outlier-clipping damage without requiring runtime absolute-maximum register scans. This validation demonstrates that the mathematical principles of QASC scale seamlessly from vision backbones to modern billion-parameter LLM configurations.

\subsection{Physical CPU Benchmarking \& Systems Compilation Insights}
\label{sec:physical_benchmarks}
To bridge the gap between our analytical hardware model and real microprocessor execution, we perform physical hardware micro-benchmarks on an Intel Xeon CPU. We measure the actual physical timing of a single FP32 low-rank projection layer ($B=256$, $D=192$, $r=8$) against its low-precision counterparts using PyTorch. 

Our physical benchmark reveals that a single FP32 projection layer executes in \textbf{0.0387 ms per batch}. However, running standard, uncompiled low-precision (BF16) operations on consumer/server CPUs in eager-mode PyTorch results in a slowdown (\textbf{0.1895 ms per batch}), representing a $0.25\times$ speedup. This behavior scales proportionally for larger, LLM-scale projection layers ($B=256$, $D=4096$, $r=16$), where uncompiled FP32 executes in \textbf{0.7615 ms} and eager BF16 takes \textbf{1.6326 ms}, giving a $0.47\times$ speedup. 

We conduct a detailed systems analysis to identify the root causes of this uncompiled framework slowdown:
\begin{itemize}
  \item \textbf{Eager-Mode Python Dispatch Overhead:} For tiny low-rank adapter projections ($r=8$), the actual computation time is extremely small (sub-microsecond range). In eager-mode frameworks like PyTorch, the execution latency is heavily dominated by Python-to-C++ CPython API wrapping, tensor shape validation, and kernel dispatching overheads, rather than the arithmetic execution itself.
  \item \textbf{Dynamic Memory Allocations:} In uncompiled execution, computing the low-rank chain $h A_k B_k$ requires allocating memory for the intermediate tensor $y = h A_k \in \mathbb{R}^{B \times r}$, writing it to DRAM, and then reloading it for the up-projection. These dynamic memory allocations completely bypass CPU caches and saturate memory bandwidth.
  \item \textbf{Dynamic Casting Instruction Stalls:} In standard eager execution, there are no fused low-precision operators. Dynamic activation quantization requires executing separate casting, rounding, and clipping operators sequentially, which stalls the CPU instruction pipeline due to frequent data conversions.
  \item \textbf{High-Level Compiler Constraints (torch.compile):} Even when compiling the low-rank projection chain in PyTorch using \texttt{torch.compile(mode="reduce-overhead")}, we measure a substantial physical slowdown on CPU (\textbf{0.0868 ms} for compiled FP32 vs \textbf{0.0314 ms} for eager FP32, giving a $0.36\times$ speedup, and \textbf{0.2547 ms} for compiled BF16, giving a $0.12\times$ speedup). This confirms that high-level compilers designed for heavy deep-learning tensors are heavily unoptimized for sub-microsecond low-rank adapters ($r=8$) on CPU. The standard compiler infrastructure introduces tracing, kernel-guard checks, and scheduling loops that exceed the raw computation time, highlighting why low-level, fused C++ kernels are necessary for on-device acceleration.
\end{itemize}

This finding underscores a critical systems-ML insight: physical speedups from low-bit PEFT quantization are never realized out-of-the-box in uncompiled environments. This validates our framing as a hardware-calibrated analytical simulation, as our ICS sandbox models the *ideal compiled execution capabilities* of a custom ONNX Runtime or ExecuTorch custom operator. In a compiled runtime, the dynamic activation quantization, down-projection, dynamic intermediate scale reduction, and de-quantization are fused into a single C++ custom operator. This allows the intermediate low-rank activations to be held entirely inside CPU vector registers or L1 cache, eliminating Python dispatch overheads, dynamic memory allocations, and casting stalls. 

We explicitly acknowledge this limitation and the simulation-to-hardware gap: while our analytical cost model indicates a projected 3.97$\times$ physical speedup under compiled execution, achieving these savings in physical deployments is highly contingent on the successful implementation of custom compiler optimizations. Without such low-level kernel fusion and compile-time integration (e.g., via ONNX Runtime custom C++ operators or ExecuTorch CustomOps), uncompiled high-level framework overheads will completely erase low-precision arithmetic benefits and even cause physical slowdowns on CPU. Compiling and fusing kernels is therefore a fundamental co-design requirement to unlock low-precision physical speedups on the edge, which we map out in our concrete systems compilation roadmap in Section~\ref{sec:compilation_roadmap}.
